% Listes :

\setlist[enumerate,1]{label=(\roman*)}

% Commandes :

\providecommand{\lesson}[3]{%
	\begin{docname}%
		\section*{\textbf{#2} #3}%
	\end{docname}%
	\begin{doctitle}%
		#3%
	\end{doctitle}%
	\begin{doccategories}%
		#1%
	\end{doccategories}%
}

\providecommand{\development}[3]{%
	\begin{docname}%
		\section*{#3}%
	\end{docname}%
	\begin{doctitle}%
		#3%
	\end{doctitle}%
	\begin{doccategories}%
		#1%
	\end{doccategories}%
}

\providecommand{\sheet}[3]{\development{#1}{#2}{#3}}

\providecommand{\ranking}[1]{%
	\begin{docname}%
		\section*{Agrégation externe #1}%
	\end{docname}%
}

\providecommand{\summary}[1]{%
	\begin{docsummary}%
		#1%
	\end{docsummary}%
	\newpar%
}

\providecommand{\reference}[2][]{%
	\begin{bookref}%
		\begin{bookrefshort}#1\end{bookrefshort}%
		\begin{bookrefpage}p. #2\end{bookrefpage}%
	\end{bookref}%
}

\providecommand{\dev}[1]{%
	\begin{devlink}#1\end{devlink}%
}

\providecommand{\newpar}{


}

% Environnements :

\theoremstyle{definition}
\newtheorem{theorem}{Théorème}
\newtheorem{property}[theorem]{Propriété}
\newtheorem{proposition}[theorem]{Proposition}
\newtheorem{lemma}[theorem]{Lemme}
\newtheorem{corollary}[theorem]{Corollaire}
\newtheorem{definition}[theorem]{Définition}
\newtheorem{application}[theorem]{Application}
\newtheorem{notation}[theorem]{Notation}
\newtheorem{example}[theorem]{Exemple}
\newtheorem{cexample}[theorem]{Contre-exemple}
\newtheorem{remark}[theorem]{Remarque}
\newtheorem{algorithm}[theorem]{Algorithme}

\renewenvironment{whitetabularx}[1]{\begin{tabularx}{\textwidth}{#1}}{\end{tabularx}}
\renewenvironment{xltabular}[2]{\begin{tabularx}{#1}{#2}}{\end{tabularx}}
\renewenvironment{flushright}{\begin{text-end}}{\end{text-end}}
